\scanpage{141}
\mspage{21}
it follows that $u$ and $n$ are co-prime and hence $q_1,\ldots,q_b$
do not divide $n$ which implies that they do not divide $v$ and
hence they do not divide $w$. Also any $p_i$ does not divide $u$
but does divide $v$ and hence does not divide $w$. Therefore none
of the prime factors $p_1,\ldots,p_a,q_1,\ldots,q_b$ of $m$ divide $w$
and hence \inserted{$g$}\deleted{$m$ and hence $g$} is a generator of $G$. Since $H$ is
of order $n$ and $n$ divides $v$ we have $\alpha(r)^v=1$ and hence
$\alpha(g)=\alpha(r)^{u+v}=\alpha(r)^u=h$.

3.\quad\uline{Quadratic Transforms of Regular Local Domains}.

Throughout this section $(R,M)$ will denote a regular local
domain of dimension $n>1$ with quotient field $K$.

\uline{Definition 1. A local domain} $(R_1,M_1)$ \uline{with quotient field}
$K$ \uline{will be called an immediate or a first quadratic transform of}
$R$ \uline{if there exists} $x\in M$, $x\notin M^2$ \uline{such that} $R_1$ \uline{is the quotient}
\uline{ring in} $K$ \uline{of} $R[Mx^{-1}]$ \uline{with respect to a maximal ideal in}
$R[Mx^{-1}]$ \uline{containing} $x$ where $R[Mx^{-1}]$ denotes the ring obtained
by adjoining to $R$ all the elements in $K$ of the form $yx^{-1}$
with $y$ in $M$; now $x\in M$ and $x\notin M^2$ implies that there ex-
ists a minimal basis $x=x_1,x_2,\ldots,x_n$ of $M$ and then
$R[Mx^{-1}]=R[x_2/x_1,\ldots,x_n/x_1]$, and hence by Lemma 1, $R_1$ is a
regular local domain of dimension $n$ having center $M$ in $R$. \uline{By}
\uline{induction we define an} $m$-\uline{th quadratic transform of} $R$ \uline{to be an}
\uline{immediate quadratic transform of an} $(m-1)$-\uline{st quadratic transform}
\uline{of} $R$. \uline{Also we shall call} $R$ \uline{its only} $0$-\uline{th quadratic transform}.
\uline{Finally if} $R^*$ \uline{is an} $m$-\uline{th quadratic transform of} $R$ \uline{for some non-}
\uline{negative integer} $m$ \uline{then we shall say that} $R^*$ \uline{is a quadratic}
\uline{transform of} $R$.
\mspage{140}
